\documentclass[10pt,twocolumn]{article}
\usepackage[T1]{fontenc}
\usepackage[utf8]{inputenc}
\usepackage{lmodern}
\usepackage[margin=1.8cm,top=2.4cm,bottom=2cm,columnsep=0.6cm]{geometry}
\usepackage{fancyhdr}
\usepackage{titlesec}
\usepackage{booktabs}
\usepackage{amsmath,amssymb,amsfonts}
\usepackage{microtype}
\usepackage{xcolor}
\usepackage{mdframed}
\usepackage{parskip}
\usepackage{enumitem}
\usepackage{hyperref}

\definecolor{rulered}{RGB}{180,30,30}
\definecolor{boxgray}{RGB}{245,245,245}
\definecolor{keywordgray}{RGB}{100,100,100}

\pagestyle{fancy}
\fancyhf{}
\fancyhead[L]{\textsc{\small Claude Mag}}
\fancyhead[C]{\small\textit{Machine Learning Research Digest}}
\fancyhead[R]{\small 2026-07-27}
\fancyfoot[C]{\small\thepage}
\renewcommand{\headrulewidth}{0.8pt}

\titleformat{\section}{\normalsize\bfseries\color{rulered}}{}{0pt}{}%
  [\vspace{-4pt}\color{rulered}\rule{\columnwidth}{0.4pt}]
\titlespacing{\section}{0pt}{8pt}{4pt}

\hypersetup{colorlinks=true,urlcolor=rulered,linkcolor=black}

\begin{document}
\setlength{\parindent}{0pt}
\setlength{\parskip}{4pt}

{\small\color{keywordgray}\textbf{Keywords:}
  reinforcement learning \textbullet{}
  RLVR \textbullet{}
  reasoning diversity \textbullet{}
  pass@k \textbullet{}
  LLM alignment}\par\vspace{4pt}

{\LARGE\bfseries\raggedright
  When RLVR Shrinks the Reasoning Boundary:
  Diagnosing Pass@$k$ Inversion}\par\vspace{4pt}

{\large\itshape\raggedright
  Reinforcement Learning with Verifiable Rewards Boosts One-Shot Accuracy
  While Silently Collapsing Sampling Diversity---a Failure Mode Traced
  to Boundary Prompts Where Rare Correct Trajectories Vanish Before
  They Can Be Reinforced}\par\vspace{6pt}

{\small\textbf{By} Todd Zhou}\par\vspace{2pt}

{\small\color{keywordgray}arXiv:2607.20543 \textbullet{} July 2026}
\par\vspace{2pt}\noindent{\color{rulered}\rule{\columnwidth}{0.6pt}}\vspace{6pt}

Reinforcement learning with verifiable rewards (RLVR) has become the
dominant recipe for training reasoning models, and its track record at
improving pass@1 is well-documented. A new preprint reveals the other
side of the ledger: RLVR can simultaneously \emph{decrease} pass@$k$
for large $k$, meaning the trained model solves fewer distinct problems
under repeated sampling than the base model it was fine-tuned from.
The author terms this phenomenon \textbf{pass@$k$ inversion} and
traces it to a concrete failure mode involving rare correct trajectories.

\begin{mdframed}[backgroundcolor=boxgray,linecolor=rulered,linewidth=1pt,
    innerleftmargin=6pt,innerrightmargin=6pt,
    innertopmargin=6pt,innerbottommargin=6pt]
\textbf{\small AT A GLANCE}\par\vspace{4pt}
\begin{description}[leftmargin=0pt,labelsep=4pt,itemsep=2pt,parsep=0pt,
    font=\small\bfseries]
  \item[Problem:] RLVR is evaluated almost exclusively at pass@1,
    masking the fact that it can reduce coverage under diverse sampling.
  \item[Why it matters:] Any system using best-of-$N$, voting, or
    repeated sampling to boost reliability may be \emph{harmed} by
    RLVR training on boundary prompts.
  \item[Method:] Pass@$k$ evaluated pre- and post-RLVR across MATH
    and GSM8K; boundary prompts identified by base model entropy;
    causal rollout analysis tracing the loss of rare correct trajectories.
  \item[Result:] RLVR improves pass@1 by 5--12\,pp while reducing
    pass@16 by 3--8\,pp on 15--25\% of test problems.
  \item[Evidence:] Per-Problem Base Anchoring (PBA) recovers 60\%
    of the pass@$k$ loss while preserving 95\% of the pass@1 gains.
\end{description}
\end{mdframed}

\section{The Big Picture}

Pass@$k$ is a standard metric for reasoning benchmarks: it asks what
fraction of problems a model can solve if given $k$ independent
sampling attempts, and a correct answer in any one attempt counts.
For a fixed pass@1 score, higher pass@$k$ reflects richer exploration
of the solution space---a property valuable when models are used
with sampling-based inference strategies such as best-of-$N$ selection,
majority voting, or Monte Carlo Tree Search.

The intuition is that if RLVR improves a model, it should improve both
pass@1 and pass@$k$. This paper shows the intuition is wrong for a
non-trivial fraction of problems. After RLVR, the policy is sharper---
it assigns higher probability to a narrower set of trajectories---and
this sharpening can eliminate rare-but-correct paths that the base
model would occasionally find at large $k$.

\section{Why Existing Approaches Fall Short}

Standard RLVR training (e.g., GRPO) collects $G$ rollouts per problem,
assigns binary rewards, and updates the policy to increase the relative
likelihood of correct rollouts over incorrect ones. This mechanism has
an implicit vulnerability: if none of the $G$ rollouts for a problem
are correct, the update treats the problem as definitively unsolvable
and suppresses all trajectories toward its vicinity in policy space.

For \textbf{boundary prompts}---problems where the base model's
correct-trajectory probability is, say, 5--15\%---a rollout group of
size $G = 8$ will observe zero correct rollouts in a large fraction
of training encounters. Each such encounter applies a gradient that
reduces the probability of rare correct trajectories, eventually
driving them to near zero even though they were genuinely present.

\section{Core Mechanism}

The author proposes a \emph{two-mode account} of post-RLVR policy
distributions:

\begin{itemize}[itemsep=2pt,parsep=0pt]
  \item \textbf{Safe prompts:} Base model has high enough
    correct-trajectory probability that at least one correct rollout
    appears reliably in each group. RLVR reinforces and concentrates
    this probability. Pass@$k$ improves monotonically with pass@1.
  \item \textbf{Boundary prompts:} Correct trajectories are too sparse
    to survive finite rollout groups. RLVR receives an absence-of-evidence
    signal and drives correct-trajectory probability toward zero.
    Pass@$k$ \emph{collapses} even as pass@1 may improve on adjacent
    safe prompts in the same batch.
\end{itemize}

Formally, for a boundary prompt $x$ with base policy correct
probability $\pi_\text{base}(x) \in [\epsilon, \delta]$ (small
$\epsilon, \delta > 0$), the expected number of correct rollouts in a
group of size $G$ is $G \cdot \pi_\text{base}(x)$. When this falls
below 1, the probability of observing zero correct rollouts in the
group is approximately $(1-\pi_\text{base})^G$, which for
$\pi_\text{base} = 0.08$ and $G = 8$ is 50.6\%.

\section{Technical Formulation}

Define pass@$k$ for a single problem $x$ and policy $\pi$ as:

\[
\mathrm{pass@}k(x, \pi)
  = 1 - \bigl(1 - \pi_\mathrm{correct}(x)\bigr)^k
\]

where $\pi_\mathrm{correct}(x) = \sum_{a \in \mathcal{A}(x)} \pi(a \mid x)$
is the policy's total mass on correct answers. Pass@$k$ inversion on
problem $x$ occurs when:

\[
\pi_\mathrm{RL,correct}(x) < \pi_\mathrm{base,correct}(x)
\]

even if $\pi_\mathrm{RL,correct}(x) > \pi_\mathrm{base,correct}(x)$
holds \emph{on average} across all problems.

The GRPO update applied to a rollout group $\{r_1, \ldots, r_G\}$
where all $r_i$ are incorrect is:

\[
\nabla \mathcal{L}_\mathrm{GRPO}
  \propto -\sum_{i=1}^G \nabla \log \pi_\theta(r_i \mid x)
\]

which suppresses the probability of all sampled trajectories. Since
rare correct trajectories were not sampled, they are inadvertently
suppressed via policy smoothing.

\section{Learning or Inference Procedure}

The Per-Problem Base Anchoring (PBA) fix proceeds as follows:

\begin{enumerate}[itemsep=2pt,parsep=0pt]
  \item Classify each training problem as boundary or safe using the
    base model's entropy $H(\pi_\mathrm{base}(\cdot \mid x))$;
    problems above a threshold $H^*$ are classified as boundary.
  \item For safe problems, apply standard RLVR with no modification.
  \item For boundary problems, augment the rollout group with
    forced correct samples from the frozen base model before computing
    the RLVR gradient.
  \item Additionally, apply a KL penalty on boundary problems:
    \[
    \mathcal{L}_\mathrm{PBA} = \mathcal{L}_\mathrm{RLVR}
      + \lambda \, D_\mathrm{KL}\!\bigl(\pi_\mathrm{base} \,\|\, \pi_\theta\bigr)
    \]
    which anchors the policy toward the base distribution on these inputs.
\end{enumerate}

\section{What the Guarantee Says}

PBA does not guarantee recovery of all lost pass@$k$ on boundary prompts.
Rather, it guarantees that the policy cannot reduce correct-trajectory
probability below a KL-bounded neighborhood of the base policy. For
$\lambda > 0$ and $D_\mathrm{KL}(\pi_\mathrm{base} \| \pi_\theta) \leq \epsilon_\lambda$,
the trained policy satisfies:

\[
\pi_\mathrm{RL,correct}(x)
  \geq \pi_\mathrm{base,correct}(x) \cdot e^{-\epsilon_\lambda}
\]

providing a fractional lower bound on retained diversity.

\section{Experimental Findings}

Experiments use Llama-3 8B and Qwen2.5 7B as base models, with
GRPO-style RLVR on MATH Level 3--5 and GSM8K problems.

\begin{center}
\small
\begin{tabular}{llcc}
\toprule
Model & Metric & Base & +RLVR \\
\midrule
Qwen2.5-7B & pass@1 & 56.3 & 68.1 \\
Qwen2.5-7B & pass@16 & 83.2 & 78.4 \\
Llama-3-8B & pass@1 & 48.9 & 57.3 \\
Llama-3-8B & pass@16 & 79.1 & 73.6 \\
\bottomrule
\end{tabular}
\end{center}

The pass@16 regressions are concentrated on boundary prompts
(15--25\% of the test set) and are negligible on safe prompts.
PBA recovers 60\% of the pass@$k$ regression on boundary prompts
while retaining 95\% of the pass@1 gain.

\section{Ablations and Interpretation}

\textbf{Group size $G$:} Larger rollout groups ($G = 16$ vs.\ $G = 8$)
reduce but do not eliminate inversion, as boundary prompts can still
yield zero correct rollouts when $\pi_\mathrm{base,correct}(x) < 1/G$.

\textbf{Problem difficulty:} Pass@$k$ inversion is uncorrelated with
pass@1 improvement magnitude; problems that yield large pass@1 gains
may or may not be on boundary prompts.

\textbf{Anchor weight $\lambda$:} Higher $\lambda$ recovers more
pass@$k$ but reduces pass@1 gains. The optimal $\lambda$ is
task-dependent.

\textbf{Broader implication:} Evaluating reasoning models solely at
pass@1 is insufficient for any deployment context using sampling-based
inference. The community should adopt pass@$k$ evaluation as standard
alongside pass@1 when reporting RLVR results.

\section*{Reference}

Todd Zhou.
\textbf{When RLVR Shrinks the Reasoning Boundary: Diagnosing
Pass@$k$ Inversion}.
arXiv:2607.20543, July 2026.
\url{https://arxiv.org/abs/2607.20543}

\end{document}
